\subNsection{1}{Reguläre Sprachen}{Kristian Popov}

Das Thema \emph{Reguläre Sprachen} unterstützt die Arbeit mit regulären Sprachen
und rechtslinearen Grammatiken.
Es bietet drei miteinander verknüpfte Funktionen:
Die Aufzählung der von einer Grammatik erzeugten Wörter, die Erzeugung einer
passenden Grammatik aus einer Sprachbeschreibung sowie die Prüfung einzelner Wörter.
Sprachbeschreibungen können entweder als Aufzählung oder in Mengenschreibweise
eingegeben werden.
Bei der Aufzählung einer Sprache werden die Wörter nach ihrer Länge sortiert.
Die Ausgabe ist dabei auf eine maximale Wortlänge und Wortanzahl begrenzt.
Bei der Wortprüfung zeigt das Thema neben der Annahme oder Ablehnung auch den Grund
für ein abweichendes Ergebnis an, beispielsweise einen nicht möglichen Übergang.
Beide Umwandlungsrichtungen sind zusätzlich als Übungen mit automatischer Rückmeldung verfügbar.

\subNsection{2}{Entscheidungen in der Architektur}{Kristian Popov}

Der Simulator besteht aus einer koordinierenden Komponente und fachlich getrennten Bausteinen.
Die Komponente \texttt{RLSimulator} verwaltet die drei Arbeitsmodi,
den Wechsel zwischen ihnen und den Einstieg in Übungen.
Sie enthält selbst keine formale Logik.
Diese liegt vollständig im Verzeichnis \texttt{domain} und ist damit
unabhängig von React testbar.

Die fachlichen Bausteine sind wie folgt aufgeteilt:

\begin{description}[leftmargin=*,style=nextline,itemsep=0.75ex]
	\item[Grammatikverarbeitung]
		Der Grammatikparser und die Validierung lesen die Eingabe ein, prüfen die
		Struktur einer rechtslinearen Grammatik und überführen sie in gemeinsame
		Datentypen.

	\item[Sprachbeschreibung]
		Ein eigener Parser verarbeitet Aufzählungen und Mengenschreibweisen.
		Er liefert eine strukturierte Beschreibung aus Faktoren, Exponenten und
		Nebenbedingungen, die anschließend in eine Grammatik umgewandelt werden
		kann.

	\item[Umwandlung und Vereinfachung]
		Die Umwandlung erzeugt aus einer Sprachbeschreibung eine strikt
		rechtslineare Grammatik.
		Ergänzende Funktionen fassen gleichwertige Nichtterminale zusammen und
		formatieren das Ergebnis für die Ausgabe.

	\item[Sprachauswertung]
		Die Aufzählung erzeugt Wörter bis zu den festgelegten Grenzen; die
		Worterkennung prüft ein einzelnes Wort gegen eine Grammatik.
		Die Übungsprüfung baut auf denselben Funktionen auf und ergänzt sie um die
		Auswertung der vorgegebenen Testfälle.
\end{description}

Wiederverwendet werden der Grammatikparser, die Grammatikvalidierung und die gemeinsamen
Datentypen.
Eine eigene Ableitungslogik direkt auf den Produktionen hätte zwei
Implementierungen erfordert, die dasselbe Verhalten abbilden müssen.
Mit der gemeinsamen Grundlage stimmen Aufzählung, Worterkennung und Übungsprüfung
stets überein.

Für die Sprachbeschreibung habe ich einen eigenen Parser geschrieben, der zwei
Formen versteht:

\begin{description}[leftmargin=*,style=nextline,itemsep=0.75ex]
	\item[Aufzählung]
		listet die Wörter unmittelbar auf:\\
		\verb|L = {"010", "0010"}|

	\item[Mengenschreibweise]
		beschreibt ein Muster mit Variablen und Nebenbedingungen:\\
		\verb+L = {a^n c b^m | n, m > 0}+
\end{description}

Endet eine Aufzählung mit einer Fortsetzungsmarkierung, versucht der Parser, daraus
ein wiederkehrendes Muster abzuleiten, statt die Wörter nur wörtlich zu übernehmen.
So lässt sich eine unendliche Sprache angeben, ohne die Mengenschreibweise schon
beherrschen zu müssen.

Bestimmte Eingaben werden bewusst abgelehnt, auch wenn sie sich formal beschreiben lassen.
Verwendet dieselbe Variable mehrere Faktoren, wie in \verb|a^n b^n|, koppelt das
deren Längen aneinander.
Der Musterfall \verb|a^n b^n| ist nicht regulär; die unterstützte Sprachbeschreibung
lässt solche gekoppelten Längenbedingungen deshalb generell nicht zu und weist sie
mit einer eigenen Meldung zurück.
Obere Schranken werden ebenfalls nicht unterstützt,
obwohl sie endliche und damit reguläre Sprachen beschreiben.
Sie würden die Eingabesprache unnötig erweitern und lassen sich klarer als Aufzählung angeben.
Diese Fälle ausdrücklich zu erklären, ist lehrreicher, als sie stillschweigend zu verwerfen.

Die erzeugte Grammatik ist strikt rechtslinear, es entstehen also ausschließlich
Regeln der Form \verb|A -> aB| und \verb|A -> a| sowie die Regel
$A \rightarrow \varepsilon$.
Eine frühere Fassung gab eine kompaktere Form wie \verb|S -> A c B| aus.
Diese Schreibweise ist zwar kürzer, aber eine allgemeine kontextfreie Regel und
verlässt damit die Typ-3-Grammatiken, um die es in diesem Thema gerade geht.

Um trotzdem eine lesbare Ausgabe zu bieten, gibt es zwei Sichten auf dasselbe
Ergebnis statt zweier Formalismen:
Die detaillierte Sicht zeigt ein Nichtterminal je Präfix, die vereinfachte Sicht
führt Nichtterminale zusammen, die dieselben Wörter ableiten.
Bei der vereinfachten Sicht werden äquivalente Teilbäume zusammengeführt.
Sie bleibt rechtslinear und kann die Grammatik nur verkleinern, nie ihre Sprache verändern.

Alle Stellen, die eine rechtslineare Grammatik ausgeben, benutzen dieselbe
Namensvergabe und Ausgabefunktion.
Zuvor lag diese Logik in drei leicht unterschiedlichen Kopien vor.
In einer davon verbrauchte ein weggelassener Endzustand einen Buchstaben, sodass in der Ausgabe
eine Lücke entstand und etwa \verb|E| übersprungen wurde.
Eine gemeinsame Funktion schließt solche Abweichungen aus und stellt sicher, dass
eine exportierte Grammatik nach dem Import dieselben Namen trägt.

\subNsection{2}{Entscheidungen im UI}{Kristian Popov}

Die drei Arbeitsweisen sind über eine Modusauswahl erreichbar.
Ihre Komponenten bleiben im Hintergrund erhalten und werden nur ein- oder ausgeblendet.
Dadurch bleiben Eingaben beim Wechsel zwischen den Modi bestehen.
Würden die Modi jeweils neu aufgebaut, ginge die begonnene Arbeit verloren
und die Ansicht spränge auf ihr Standardbeispiel zurück.

Ist eine Übung geöffnet, ersetzt eine geführte Arbeitsfläche den freien Umwandler.
Vorgabe und Aufgabe stehen dabei nebeneinander:
In der einen Richtung ist die Grammatik gegeben und die Sprache gesucht, in der
anderen umgekehrt.
Ein Moduswechsel während einer offenen Übung würde den aktuellen Bearbeitungsstand
verwerfen. Deshalb wird vorher eine Bestätigung verlangt.

Fehlermeldungen erscheinen unmittelbar unter dem jeweiligen Eingabefeld, aber erst,
nachdem überhaupt etwas eingegeben wurde.
Ein leeres Feld ist kein Fehler, sondern der Ausgangszustand.
Zusätzlich lässt sich die unterstützte Syntax direkt neben dem Eingabefeld
aufklappen, damit die Schreibweise nicht erst gesucht werden muss.

\subNsection{2}{Andere Entscheidungen}{Kristian Popov}

Die Kontrolle der Übungen vergleicht keine Texte.
Für die Richtung \emph{Sprache}$\rightarrow$ \emph{Grammatik} sind zu jeder Aufgabe Wörter hinterlegt,
die erzeugt werden müssen, und Wörter, die nicht erzeugt werden dürfen.
Geprüft wird, ob die eingegebene Grammatik alle geforderten Wörter ableitet und
alle ausgeschlossenen ablehnt.
Ein Textvergleich mit einer Musterlösung wäre hier unbrauchbar, da dieselbe Sprache
von vielen unterschiedlichen Grammatiken erzeugt werden kann und die Sprachen zum
Teil unendlich sind.
Die Rückmeldung nennt daher falsch erzeugte und fehlende Wörter statt nur einer
Ja-Nein-Aussage.
Da nur vorbereitete Beispiele geprüft werden, dient auch diese Kontrolle
der gezielten Rückmeldung und ist kein allgemeiner Beweis für die Gleichheit zweier Sprachen.

Für die Richtung \emph{Grammatik} $\rightarrow$ \emph{Sprache} verwendet die
Kontrolle dieselben Grenzen für Wortlänge und Wortanzahl wie die Schaltfläche zum
Anzeigen der Lösung.
Andernfalls könnte eine angezeigte Lösung die anschließende Prüfung nicht bestehen.

Beim Import und Export unterscheiden sich die beiden Formate bewusst.
Eine Grammatik wird strukturiert kodiert, also in Nichtterminale, Terminale und
Produktionen zerlegt und bitweise gepackt.
Eine Sprachbeschreibung wird dagegen als Text übernommen, weil sie vom Benutzer
selbst formuliert ist und ihre Schreibweise erhalten bleiben soll.
Beide Seiten prüfen beim Kodieren und beim Dekodieren, dass der Inhalt gültig ist.

Jeder Modus speichert seinen Stand in einem eigenen Bereich der lokalen Ablage.
Ohne diese Trennung hätten sich die Grammatik des Aufzählungsmodus und die des
Worterkenners gegenseitig überschrieben, da beide auf demselben Eingabetyp arbeiten.

Die Aufgabenstellungen selbst sind in einer sprachneutralen Notation hinterlegt und
nicht in den Übersetzungsdateien.
Eine Produktion wie \verb|START = "1", A| ist in beiden Oberflächensprachen
identisch; eine Übersetzung würde hier nur zusätzliche Fehlerquellen schaffen.
